% ==================== Page and typography ====================
\usepackage[a4paper,margin=2.5cm,headheight=15pt]{geometry}
\usepackage{fontspec}
\usepackage{microtype}
\usepackage{setspace}
\usepackage{fancyhdr}
\usepackage{tocloft}
\usepackage{xcolor}

\definecolor{noteblue}{HTML}{245B8A}
\definecolor{notegreen}{HTML}{2D6A4F}
\definecolor{notered}{HTML}{A23E48}
\definecolor{notegrey}{HTML}{F4F6F8}

\setstretch{1.18}
\setlength{\parindent}{2em}
\setlength{\parskip}{0.25em}
\pagestyle{fancy}
\fancyhf{}
\fancyhead[LE]{\small\leftmark}
\fancyhead[RO]{\small\rightmark}
\fancyfoot[C]{\thepage}

% Use leaders and right-aligned page numbers for chapter entries in the TOC.
\renewcommand{\cftchapleader}{\cftdotfill{\cftdotsep}}

% ==================== Mathematics and physics ====================
\usepackage{amsmath,amssymb,amsthm}
\usepackage{mathtools}
\usepackage{bm}
\usepackage{tensor}
\usepackage{braket}
\usepackage{cancel}
\usepackage{esint}
\usepackage{empheq}

% Units, numerical values, and chemical/nuclear equations
\usepackage{siunitx}
\sisetup{
  separate-uncertainty = true,
  per-mode = symbol,
  range-phrase = {--},
  range-units = single
}
\usepackage[version=4]{mhchem}

% ==================== Figures, tables, and diagrams ====================
\usepackage{graphicx}
\graphicspath{{images/}{../images/}}
\usepackage{float}
\usepackage{subcaption}
\usepackage{booktabs}
\usepackage{array}
\usepackage{longtable}
\usepackage{multirow}
\usepackage{tabularx}

\usepackage{tikz}
\usetikzlibrary{
  arrows.meta,
  angles,
  quotes,
  calc,
  decorations.pathmorphing,
  patterns,
  positioning
}
\usepackage{pgfplots}
\pgfplotsset{compat=1.18}
\usepackage[europeanresistors,americaninductors]{circuitikz}

% Optional: Feynman diagrams. LuaLaTeX is recommended when enabled.
% \usepackage{tikz-feynman}

% ==================== Code, lists, and note boxes ====================
\usepackage{enumitem}
\usepackage{listings}
\usepackage[most]{tcolorbox}

\lstset{
  basicstyle=\ttfamily\small,
  numbers=left,
  numberstyle=\tiny\color{gray},
  frame=single,
  breaklines=true,
  backgroundcolor=\color{notegrey}
}

\newtcolorbox{summarybox}[1][Section Summary]{
  enhanced,
  breakable,
  colback=noteblue!5,
  colframe=noteblue,
  fonttitle=\bfseries,
  title=#1
}
\newtcolorbox{warningbox}[1][Common Pitfall]{
  enhanced,
  breakable,
  colback=notered!5,
  colframe=notered,
  fonttitle=\bfseries,
  title=#1
}

% ==================== Theorem environments ====================
\theoremstyle{definition}
\newtheorem{definition}{Definition}[chapter]
\newtheorem{example}{Example}[chapter]
\theoremstyle{plain}
\newtheorem{theorem}{Theorem}[chapter]
\newtheorem{law}{Law}[chapter]
\newtheorem{proposition}{Proposition}[chapter]
\theoremstyle{remark}
\newtheorem{remark}{Remark}[chapter]

% ==================== Subfiles and bibliography ====================
\usepackage{subfiles}
\usepackage[
  backend=biber,
  style=numeric,
  sorting=none,
  giveninits=true
]{biblatex}
% \subfix automatically corrects the path when compiling chapters/*.tex.
\addbibresource{\subfix{references.bib}}

% ==================== Links and intelligent references ====================
% Load these packages near the end of the preamble.
\usepackage[
  colorlinks=true,
  linkcolor=noteblue,
  citecolor=notegreen,
  urlcolor=notered,
  bookmarksopen=true
]{hyperref}
\usepackage[nameinlink,noabbrev]{cleveref}

\crefname{equation}{equation}{equations}
\crefname{figure}{figure}{figures}
\crefname{table}{table}{tables}
\crefname{definition}{definition}{definitions}
\crefname{theorem}{theorem}{theorems}
\crefname{law}{law}{laws}
\crefname{example}{example}{examples}

% ==================== Common physics commands ====================
\newcommand{\dd}{\mathop{}\!\mathrm{d}}
\newcommand{\ee}{\mathrm{e}}
\newcommand{\ii}{\mathrm{i}}
\newcommand{\vect}[1]{\bm{#1}}
\newcommand{\uvect}[1]{\bm{\hat{#1}}}
\newcommand{\grad}{\bm{\nabla}}
\newcommand{\laplacian}{\nabla^2}
\newcommand{\comm}[2]{\left[#1,#2\right]}
\newcommand{\acomm}[2]{\left\{#1,#2\right\}}
\newcommand{\expect}[1]{\left\langle #1\right\rangle}
\newcommand{\absq}[1]{\left|#1\right|^2}
\newcommand{\evalat}[2]{\left.#1\right|_{#2}}
\NewDocumentCommand{\dv}{m m}{\frac{\dd #1}{\dd #2}}
\NewDocumentCommand{\pdv}{m m}{\frac{\partial #1}{\partial #2}}
\DeclareMathOperator{\Tr}{Tr}
\DeclareMathOperator{\Res}{Res}
\DeclareMathOperator{\sgn}{sgn}

\numberwithin{equation}{chapter}
